% GENERATED FILE. Edit canonical lessons, then run npm run generate:books.
% canonical-source-hash: fnv1a64:f6c07e379f6bd9ca
% canonical-lessons: TA-C323-iruticcatanku, TA-C323-urvalam, TA-C323-malarvalaiyam, TA-C323-putai, TA-C323-kallarai

\chapter{A Funeral}
\label{ch:ta-a-funeral}
\hlchaptermodality{323}

\begin{chapteropening}
\textbf{By the end of this chapter:} \emph{I can talk about a funeral with respect.}

Complete the last lesson of chapter 323: I can talk about a funeral with respect.
\end{chapteropening}

\section[\texorpdfstring{\ta{இறுதிச்சடங்கு} (iṟuticcaṭaṅku)}{இறுதிச்சடங்கு (iṟuticcaṭaṅku)}]{\ta{இறுதிச்சடங்கு} (iṟuticcaṭaṅku) — a funeral}

\label{lesson:TA-C323-iruticcatanku}



\begin{quote}
Before the new one: say the Tamil for a wedding dais, then the Tamil for a pandal.
\end{quote}

\subsection*{You'll want to know: \ta{இறுதிச்சடங்கு}}
\textbf{\ta{இறுதிச்சடங்கு}} — \emph{iṟuticcaṭaṅku} — \textquotedblleft{}a funeral\textquotedblright{}.

\begin{quote}
\ta{அவருடைய} \ta{இறுதிச்சடங்கில்} \ta{ஊரே} \ta{கலந்துகொண்டது}.
\end{quote}

\begin{quote}
\emph{avaruṭaiya iṟuticcaṭaṅkil ūrē kalantukoṇṭatu.}
\end{quote}

\begin{quote}
The whole village attended his funeral.
\end{quote}

\ta{இறுதி}, end + \ta{சடங்கு}, rite.

\begin{cousinweb}
\textbf{\ta{இறுதிச்சடங்கு}} (\emph{iṟuticcaṭaṅku}) is a Tamil compound built from \ta{இறுதி}, \textquotedblleft{}end, death\textquotedblright{} (DEDR 514), and \ta{சடங்கு}, \textquotedblleft{}a rite\textquotedblright{}.
\end{cousinweb}

\begin{grammarlens}[title={What you've built}]
One more word for telling what happened.
\end{grammarlens}

\subsection*{Your turn}
\begin{itemize}
\raggedright
  \item \emph{Say it:} \emph{iṟuticcaṭaṅku}
  \item \emph{Say it:} \emph{iṟuticcaṭaṅku}, once more
  \item \emph{Say it:} \emph{avaruṭaiya iṟuticcaṭaṅkil ūrē kalantukoṇṭatu.} — The whole village attended his funeral.
  \item \emph{Recall:} say the Tamil for a wedding dais, then the Tamil for a pandal, then say \emph{iṟuticcaṭaṅku} again
\end{itemize}

\begin{tcolorbox}[breakable,colback=teal!4,colframe=teal!35!black,title={Before you move on}]
What does \ta{இறுதிச்சடங்கு} mean? (\textquotedblleft{}A funeral\textquotedblright{}.) Say it once more.
\end{tcolorbox}
\section[\texorpdfstring{\ta{ஊர்வலம்} (ūrvalam)}{ஊர்வலம் (ūrvalam)}]{\ta{ஊர்வலம்} (ūrvalam) — a procession}

\label{lesson:TA-C323-urvalam}



\begin{quote}
Before the new one: say the Tamil for a pandal, then the Tamil for a funeral.
\end{quote}

\subsection*{You'll want to know: \ta{ஊர்வலம்}}
\textbf{\ta{ஊர்வலம்}} — \emph{ūrvalam} — \textquotedblleft{}a procession\textquotedblright{}.

\begin{quote}
\ta{இறுதி} \ta{ஊர்வலம்} \ta{மெதுவாகச்} \ta{சென்றது}.
\end{quote}

\begin{quote}
\emph{iṟuti ūrvalam metuvākac ceṉṟatu.}
\end{quote}

\begin{quote}
The funeral procession moved slowly.
\end{quote}

Also for festivals and weddings: \ta{ஊர்}, town + \ta{வலம்}, going round.

\begin{cousinweb}
\textbf{\ta{ஊர்வலம்}} (\emph{ūrvalam}) is a Tamil compound built from \ta{ஊர்}, \textquotedblleft{}town\textquotedblright{} (DEDR 752), and \ta{வலம்}, \textquotedblleft{}the right side\textquotedblright{}, as in going round to the right (DEDR 5276).
\end{cousinweb}

\begin{grammarlens}[title={What you've built}]
One more word for telling what happened.
\end{grammarlens}

\subsection*{Your turn}
\begin{itemize}
\raggedright
  \item \emph{Say it:} \emph{ūrvalam}
  \item \emph{Say it:} \emph{ūrvalam}, once more
  \item \emph{Say it:} \emph{iṟuti ūrvalam metuvākac ceṉṟatu.} — The funeral procession moved slowly.
  \item \emph{Recall:} say the Tamil for a pandal, then the Tamil for a funeral, then say \emph{ūrvalam} again
\end{itemize}

\begin{tcolorbox}[breakable,colback=teal!4,colframe=teal!35!black,title={Before you move on}]
What does \ta{ஊர்வலம்} mean? (\textquotedblleft{}A procession\textquotedblright{}.) Say it once more.
\end{tcolorbox}
\section[\texorpdfstring{\ta{மலர்வளையம்} (malarvaḷaiyam)}{மலர்வளையம் (malarvaḷaiyam)}]{\ta{மலர்வளையம்} (malarvaḷaiyam) — a wreath}

\label{lesson:TA-C323-malarvalaiyam}



\begin{quote}
Before the new one: say the Tamil for a funeral, then the Tamil for a procession.
\end{quote}

\subsection*{You'll want to know: \ta{மலர்வளையம்}}
\textbf{\ta{மலர்வளையம்}} — \emph{malarvaḷaiyam} — \textquotedblleft{}a wreath\textquotedblright{}.

\begin{quote}
\ta{நண்பர்கள்} \ta{ஒரு} \ta{மலர்வளையம்} \ta{வைத்தார்கள்}.
\end{quote}

\begin{quote}
\emph{naṇparkaḷ oru malarvaḷaiyam vaittārkaḷ.}
\end{quote}

\begin{quote}
Friends laid a wreath.
\end{quote}

\ta{மலர்வளையம்} \ta{வை}, to lay a wreath.

\begin{cousinweb}
\textbf{\ta{மலர்வளையம்}} (\emph{malarvaḷaiyam}) is a Tamil compound built from \ta{மலர்}, \textquotedblleft{}flower\textquotedblright{} (DEDR 4739), and \ta{வளையம்}, \textquotedblleft{}ring, circle\textquotedblright{} (DEDR 5313).
\end{cousinweb}

\begin{grammarlens}[title={What you've built}]
One more word for telling what happened.
\end{grammarlens}

\subsection*{Your turn}
\begin{itemize}
\raggedright
  \item \emph{Say it:} \emph{malarvaḷaiyam}
  \item \emph{Say it:} \emph{malarvaḷaiyam}, once more
  \item \emph{Say it:} \emph{naṇparkaḷ oru malarvaḷaiyam vaittārkaḷ.} — Friends laid a wreath.
  \item \emph{Recall:} say the Tamil for a funeral, then the Tamil for a procession, then say \emph{malarvaḷaiyam} again
\end{itemize}

\begin{tcolorbox}[breakable,colback=teal!4,colframe=teal!35!black,title={Before you move on}]
What does \ta{மலர்வளையம்} mean? (\textquotedblleft{}A wreath\textquotedblright{}.) Say it once more.
\end{tcolorbox}
\section[\texorpdfstring{\ta{புதை} (putai)}{புதை (putai)}]{\ta{புதை} (putai) — to bury}

\label{lesson:TA-C323-putai}



\begin{quote}
Before the new one: say the Tamil for a procession, then the Tamil for a wreath.
\end{quote}

\subsection*{You'll want to know: \ta{புதை}}
\textbf{\ta{புதை}} — \emph{putai} — \textquotedblleft{}to bury\textquotedblright{}.

\begin{quote}
\ta{அவர்கள்} \ta{இறந்த} \ta{நாயைத்} \ta{தோட்டத்தில்} \ta{புதைத்தார்கள்}.
\end{quote}

\begin{quote}
\emph{avarkaḷ iṟanta nāyait tōṭṭattil putaittārkaḷ.}
\end{quote}

\begin{quote}
They buried the dead dog in the garden.
\end{quote}

Past (I …): \ta{புதைத்தேன்}, \emph{putaittēṉ}.

Without an object: \ta{புதைந்தது}, it was buried or sank in.

\begin{cousinweb}
\textbf{\ta{புதை}} (\emph{putai}) is native Dravidian (DEDR 4509): to be buried, hidden, and with -\ta{த்த}- to bury.
\end{cousinweb}

\begin{grammarlens}[title={What you've built}]
One more word for telling what happened.
\end{grammarlens}

\subsection*{Your turn}
\begin{itemize}
\raggedright
  \item \emph{Say it:} \emph{putai}
  \item \emph{Say it:} \emph{putai}, once more
  \item \emph{Say it:} \emph{avarkaḷ iṟanta nāyait tōṭṭattil putaittārkaḷ.} — They buried the dead dog in the garden.
  \item \emph{Recall:} say the Tamil for a procession, then the Tamil for a wreath, then say \emph{putai} again
\end{itemize}

\begin{tcolorbox}[breakable,colback=teal!4,colframe=teal!35!black,title={Before you move on}]
What does \ta{புதை} mean? (\textquotedblleft{}To bury\textquotedblright{}.) Say it once more.
\end{tcolorbox}
\section[\texorpdfstring{\ta{கல்லறை} (kallaṟai)}{கல்லறை (kallaṟai)}]{\ta{கல்லறை} (kallaṟai) — a grave; a tomb}

\label{lesson:TA-C323-kallarai}



\begin{quote}
Before the new one: say the Tamil for a wreath, then the Tamil for to bury.
\end{quote}

\subsection*{You'll want to know: \ta{கல்லறை}}
\textbf{\ta{கல்லறை}} — \emph{kallaṟai} — \textquotedblleft{}a grave; a tomb\textquotedblright{}.

\begin{quote}
\ta{அவருடைய} \ta{கல்லறை} \ta{ஊருக்கு} \ta{வெளியே} \ta{இருக்கிறது}.
\end{quote}

\begin{quote}
\emph{avaruṭaiya kallaṟai ūrukku veḷiyē irukkiṟatu.}
\end{quote}

\begin{quote}
His grave is outside the village.
\end{quote}

\ta{கல்}, stone + \ta{அறை}, chamber.

\begin{cousinweb}
\textbf{\ta{கல்லறை}} (\emph{kallaṟai}) is a Tamil compound built from \ta{கல்}, \textquotedblleft{}stone\textquotedblright{} (DEDR 1298), and \ta{அறை}, \textquotedblleft{}a room, a chamber\textquotedblright{}.
\end{cousinweb}

\begin{grammarlens}[title={What you've built}]
One more word for telling what happened.
\end{grammarlens}

\subsection*{Your turn}
\begin{itemize}
\raggedright
  \item \emph{Say it:} \emph{kallaṟai}
  \item \emph{Say it:} \emph{kallaṟai}, once more
  \item \emph{Say it:} \emph{avaruṭaiya kallaṟai ūrukku veḷiyē irukkiṟatu.} — His grave is outside the village.
  \item \emph{Recall:} say the Tamil for a wreath, then the Tamil for to bury, then say \emph{kallaṟai} again
\end{itemize}

\begin{tcolorbox}[breakable,colback=teal!4,colframe=teal!35!black,title={Before you move on}]
What does \ta{கல்லறை} mean? (\textquotedblleft{}A grave; a tomb\textquotedblright{}.) Say it once more.
\end{tcolorbox}
